% GENERATED FILE. Edit canonical lessons, then run npm run generate:books.
% canonical-source-hash: fnv1a64:6d5e9a9850bd1749
% canonical-lessons: BN-C357-probondho, BN-C357-bishoy, BN-C357-montobyo, BN-C357-uposonghar, BN-C357-dharona

\chapter{Reporting Findings}
\label{ch:bn-reporting-findings}
\hlchaptermodality{357}

\begin{chapteropening}
\textbf{By the end of this chapter:} \emph{I can say what a piece of research is about and how it ends.}

Complete the last lesson of chapter 357: I can say what a piece of research is about and how it ends.
\end{chapteropening}

\section[\texorpdfstring{\bn{প্রবন্ধ} (prôbôndhô)}{প্রবন্ধ (prôbôndhô)}]{\bn{প্রবন্ধ} (prôbôndhô) — an essay; an article}

\label{lesson:BN-C357-probondho}



\begin{quote}
Before the new one: say the Bengali for a yardstick, then the Bengali for to assess.
\end{quote}

\subsection*{You'll want to know: \bn{প্রবন্ধ}}
\textbf{\bn{প্রবন্ধ}} — \emph{prôbôndhô} — \textquotedblleft{}an essay; an article\textquotedblright{}.

\begin{quote}
\bn{মীরা} \bn{পত্রিকায়} \bn{একটা} \bn{প্রবন্ধ} \bn{লিখল।}
\end{quote}

\begin{quote}
\emph{mīrā pôtrikāy ekṭā prôbôndhô likhlô.}
\end{quote}

\begin{quote}
Mira wrote an article in a magazine.
\end{quote}

\begin{cousinweb}
\textbf{\bn{প্রবন্ধ}} (\emph{prôbôndhô}) is Sanskrit \emph{prabandha}, a binding together (Monier-Williams), and so a connected piece of writing. Bengali uses it for an essay or an article.
\end{cousinweb}

\begin{grammarlens}[title={What you've built}]
One more word for giving reasons and opinions.
\end{grammarlens}

\subsection*{Your turn}
\begin{itemize}
\raggedright
  \item \emph{Say it:} \emph{prôbôndhô}
  \item \emph{Say it:} \emph{prôbôndhô}, once more
  \item \emph{Say it:} \emph{mīrā pôtrikāy ekṭā prôbôndhô likhlô.} — Mira wrote an article in a magazine.
  \item \emph{Recall:} say the Bengali for a yardstick, then the Bengali for to assess, then say \emph{prôbôndhô} again
\end{itemize}

\begin{tcolorbox}[breakable,colback=teal!4,colframe=teal!35!black,title={Before you move on}]
What does \bn{প্রবন্ধ} mean? (\textquotedblleft{}An essay; an article\textquotedblright{}.) Say it once more.
\end{tcolorbox}
\section[\texorpdfstring{\bn{বিষয়} (bishôy)}{বিষয় (bishôy)}]{\bn{বিষয়} (bishôy) — a subject; a topic}

\label{lesson:BN-C357-bishoy}



\begin{quote}
Before the new one: say the Bengali for to assess, then the Bengali for an essay.
\end{quote}

\subsection*{You'll want to know: \bn{বিষয়}}
\textbf{\bn{বিষয়}} — \emph{bishôy} — \textquotedblleft{}a subject; a topic\textquotedblright{}.

\begin{quote}
\bn{তোমার} \bn{প্রবন্ধের} \bn{বিষয়} \bn{কী}?
\end{quote}

\begin{quote}
\emph{tomār prôbôndher bishôy kī?}
\end{quote}

\begin{quote}
What is your essay about?
\end{quote}

\begin{cousinweb}
\textbf{\bn{বিষয়}} (\emph{bishôy}) is Sanskrit \emph{viṣaya}, a sphere or an object of attention (Monier-Williams). In Bengali it is the subject of a talk, a book or a school course.
\end{cousinweb}

\begin{grammarlens}[title={What you've built}]
One more word for giving reasons and opinions.
\end{grammarlens}

\subsection*{Your turn}
\begin{itemize}
\raggedright
  \item \emph{Say it:} \emph{bishôy}
  \item \emph{Say it:} \emph{bishôy}, once more
  \item \emph{Say it:} \emph{tomār prôbôndher bishôy kī?} — What is your essay about?
  \item \emph{Recall:} say the Bengali for to assess, then the Bengali for an essay, then say \emph{bishôy} again
\end{itemize}

\begin{tcolorbox}[breakable,colback=teal!4,colframe=teal!35!black,title={Before you move on}]
What does \bn{বিষয়} mean? (\textquotedblleft{}A subject; a topic\textquotedblright{}.) Say it once more.
\end{tcolorbox}
\section[\texorpdfstring{\bn{মন্তব্য} (môntôbyô)}{মন্তব্য (môntôbyô)}]{\bn{মন্তব্য} (môntôbyô) — a comment; a remark}

\label{lesson:BN-C357-montobyo}



\begin{quote}
Before the new one: say the Bengali for an essay, then the Bengali for a subject.
\end{quote}

\subsection*{You'll want to know: \bn{মন্তব্য}}
\textbf{\bn{মন্তব্য}} — \emph{môntôbyô} — \textquotedblleft{}a comment; a remark\textquotedblright{}.

\begin{quote}
\bn{তোমার} \bn{মন্তব্যটা} \bn{দরকারি} \bn{ছিল।}
\end{quote}

\begin{quote}
\emph{tomār môntôbyôṭā dôrkāri chhilô.}
\end{quote}

\begin{quote}
Your comment was useful.
\end{quote}

\begin{cousinweb}
\textbf{\bn{মন্তব্য}} (\emph{môntôbyô}) is Sanskrit \emph{mantavya}, 'what is to be thought' (Monier-Williams). Bengali uses it for a short remark or comment.
\end{cousinweb}

\begin{grammarlens}[title={What you've built}]
One more word for giving reasons and opinions.
\end{grammarlens}

\subsection*{Your turn}
\begin{itemize}
\raggedright
  \item \emph{Say it:} \emph{môntôbyô}
  \item \emph{Say it:} \emph{môntôbyô}, once more
  \item \emph{Say it:} \emph{tomār môntôbyôṭā dôrkāri chhilô.} — Your comment was useful.
  \item \emph{Recall:} say the Bengali for an essay, then the Bengali for a subject, then say \emph{môntôbyô} again
\end{itemize}

\begin{tcolorbox}[breakable,colback=teal!4,colframe=teal!35!black,title={Before you move on}]
What does \bn{মন্তব্য} mean? (\textquotedblleft{}A comment; a remark\textquotedblright{}.) Say it once more.
\end{tcolorbox}
\section[\texorpdfstring{\bn{উপসংহার} (upôsônghār)}{উপসংহার (upôsônghār)}]{\bn{উপসংহার} (upôsônghār) — a conclusion of a text; a closing section}

\label{lesson:BN-C357-uposonghar}



\begin{quote}
Before the new one: say the Bengali for a subject, then the Bengali for a remark.
\end{quote}

\subsection*{You'll want to know: \bn{উপসংহার}}
\textbf{\bn{উপসংহার}} — \emph{upôsônghār} — \textquotedblleft{}a conclusion of a text; a closing section\textquotedblright{}.

\begin{quote}
\bn{প্রবন্ধের} \bn{উপসংহারটা} \bn{ছোট।}
\end{quote}

\begin{quote}
\emph{prôbôndher upôsônghārṭā chhoṭo.}
\end{quote}

\begin{quote}
The essay's conclusion is short.
\end{quote}

\begin{cousinweb}
\textbf{\bn{উপসংহার}} (\emph{upôsônghār}) is Sanskrit \emph{upasaṃhāra}, summing up (Monier-Williams): the closing part of an essay or a speech.
\end{cousinweb}

\begin{grammarlens}[title={What you've built}]
One more word for giving reasons and opinions.
\end{grammarlens}

\subsection*{Your turn}
\begin{itemize}
\raggedright
  \item \emph{Say it:} \emph{upôsônghār}
  \item \emph{Say it:} \emph{upôsônghār}, once more
  \item \emph{Say it:} \emph{prôbôndher upôsônghārṭā chhoṭo.} — The essay's conclusion is short.
  \item \emph{Recall:} say the Bengali for a subject, then the Bengali for a remark, then say \emph{upôsônghār} again
\end{itemize}

\begin{tcolorbox}[breakable,colback=teal!4,colframe=teal!35!black,title={Before you move on}]
What does \bn{উপসংহার} mean? (\textquotedblleft{}A conclusion of a text; a closing section\textquotedblright{}.) Say it once more.
\end{tcolorbox}
\section[\texorpdfstring{\bn{ধারণা} (dhārônā)}{ধারণা (dhārônā)}]{\bn{ধারণা} (dhārônā) — an idea; a notion}

\label{lesson:BN-C357-dharona}



\begin{quote}
Before the new one: say the Bengali for a remark, then the Bengali for a conclusion of a text.
\end{quote}

\subsection*{You'll want to know: \bn{ধারণা}}
\textbf{\bn{ধারণা}} — \emph{dhārônā} — \textquotedblleft{}an idea; a notion\textquotedblright{}.

\begin{quote}
\bn{এই} \bn{বিষয়ে} \bn{আমার} \bn{কোনো} \bn{ধারণা} \bn{নেই।}
\end{quote}

\begin{quote}
\emph{ei bishôye āmār kono dhārônā nei.}
\end{quote}

\begin{quote}
I have no idea about this subject.
\end{quote}

\begin{cousinweb}
\textbf{\bn{ধারণা}} (\emph{dhārônā}) is Sanskrit \emph{dhāraṇā}, holding, keeping in mind (Monier-Williams): what you hold in your head about something, a notion.
\end{cousinweb}

\begin{grammarlens}[title={What you've built}]
One more word for giving reasons and opinions.
\end{grammarlens}

\subsection*{Your turn}
\begin{itemize}
\raggedright
  \item \emph{Say it:} \emph{dhārônā}
  \item \emph{Say it:} \emph{dhārônā}, once more
  \item \emph{Say it:} \emph{ei bishôye āmār kono dhārônā nei.} — I have no idea about this subject.
  \item \emph{Recall:} say the Bengali for a remark, then the Bengali for a conclusion of a text, then say \emph{dhārônā} again
\end{itemize}

\begin{tcolorbox}[breakable,colback=teal!4,colframe=teal!35!black,title={Before you move on}]
What does \bn{ধারণা} mean? (\textquotedblleft{}An idea; a notion\textquotedblright{}.) Say it once more.
\end{tcolorbox}
